\chapter{La piattaforma FrameSAI}

La fase di sviluppo ha tradotto il prototipo ad alta fedeltà in una piattaforma web navigabile, mantenendo la struttura concettuale definita durante la progettazione e introducendo i meccanismi necessari affinché le relazioni tra gli elementi del framework non rimanessero soltanto rappresentate graficamente. La versione implementata non è stata concepita come un archivio lineare di contenuti, ma come uno spazio di consultazione nel quale Principi, Linee Guida, Criteri di Successo e Design Pattern possono essere raggiunti secondo percorsi differenti. Tale impostazione deriva dalla natura stessa di FrameSAI: pur conservando una progressione dal generale al particolare, il framework contiene relazioni molti-a-molti che rendono insufficiente una navigazione esclusivamente sequenziale.

La piattaforma rende quindi operativo il modello informativo discusso nel capitolo precedente. I quattro Principi costituiscono il livello di orientamento teorico; le dodici Linee Guida traducono tali principi in obblighi riferiti al comportamento del sistema; i sedici Criteri di Successo consentono di osservare condizioni verificabili; i Design Pattern forniscono soluzioni progettuali contestualizzate. La versione corrente espone quattordici pattern per i quali la fonte di progetto contiene una scheda completa. La numerazione originaria è stata conservata, compreso il passaggio da DP8 a DP10, evitando di introdurre un contenuto DP9 non documentato. Questa scelta mantiene la corrispondenza tra piattaforma e materiale di ricerca, anche quando la sequenza numerica non risulta continua.

L'implementazione ha richiesto un adattamento rispetto alle schermate realizzate in Figma. Nel prototipo, infatti, dimensioni e posizione degli elementi erano state definite assumendo un'area di lavoro stabile; nel browser è stato invece necessario considerare larghezze variabili, scorrimento verticale, indirizzi raggiungibili direttamente e persistenza della navigazione. Il risultato conserva il linguaggio visuale progettato nella fase Hi-Fi, ma lo inserisce in una struttura responsive e verificabile. In questo passaggio, l'interfaccia non è stata trattata come una semplice trasposizione estetica: le scelte visuali sono state ricondotte a componenti riutilizzabili, mentre i contenuti e le loro associazioni sono stati separati dalla presentazione.

\section{Illustrazione percorso nel dettaglio}

Il punto di ingresso della piattaforma è la homepage, che presenta FrameSAI attraverso una breve descrizione e rende immediatamente visibili le dimensioni del framework. I valori riportati non hanno soltanto una funzione riepilogativa, ma anticipano l'estensione delle quattro librerie e permettono all'utente di costruire una prima rappresentazione dell'ambiente. Dalla stessa schermata è possibile iniziare l'esplorazione generale oppure accedere direttamente alla libreria dei Design Pattern. Questa doppia possibilità risponde a due esigenze differenti: da una parte l'avvicinamento progressivo di chi non conosce ancora il framework, dall'altra l'accesso rapido di chi è già alla ricerca di una soluzione progettuale.

La navigazione principale rimane disponibile nella parte sinistra dello schermo e contiene i collegamenti a Homepage, Principles, Guidelines, Success Criteria e Design Patterns. La voce relativa alla sezione corrente è evidenziata attraverso una regione più chiara, rendendo percepibile la posizione dell'utente senza richiedere la ricostruzione del percorso precedente. Su schermi desktop la barra è fissata al margine della finestra e non segue lo scorrimento del contenuto: anche nelle librerie più estese, le destinazioni principali restano quindi disponibili. La sua altezza e la dimensione del logo sono state calibrate affinché tutte le voci entrino nella visualizzazione senza produrre uno scorrimento interno. Su dispositivi con larghezza ridotta, la medesima struttura viene trasformata in una barra superiore; le cinque destinazioni restano contemporaneamente visibili e non viene introdotto un secondo livello di menu.

Il percorso più aderente alla logica top-down prende avvio dalla Library of Principles. Le quattro card identificano Transparency, Fairness, Automation Level e Protection attraverso codice, titolo, colore e descrizione sintetica. La selezione di una card conduce alla pagina del principio, nella quale la definizione teorica viene accompagnata da un esempio applicativo. La pagina non costituisce un punto terminale: tre collegamenti permettono di proseguire verso Linee Guida, Criteri di Successo e Design Pattern. La freccia collocata nell'intestazione riporta invece alla Library of Principles. È stata assegnata una destinazione deterministica alla freccia, anziché affidarla alla cronologia del browser, poiché la stessa pagina può essere raggiunta da percorsi differenti o tramite un indirizzo diretto. In questo modo il comando conserva sempre lo stesso significato e riconduce al livello gerarchico immediatamente superiore.

La sezione delle Linee Guida mette in evidenza il carattere trasversale delle associazioni. Ogni card mostra il codice della guideline, la formulazione completa e i Principi ai quali essa è collegata. A differenza di una lista statica, la pagina consente di filtrare le dodici Linee Guida selezionando uno dei quattro Principi. I colori rimangono coerenti con quelli utilizzati nella Library of Principles, ma sono accompagnati dalle sigle P1, P2, P3 e P4 affinché la relazione non dipenda unicamente dalla percezione cromatica. Nella parte inferiore della card sono riportati il numero di Design Pattern e di Success Criteria connessi alla Linea Guida. Tali valori sono collegamenti effettivi: la loro selezione apre la libreria corrispondente applicando il filtro appropriato. Il passaggio tra i livelli del framework avviene pertanto conservando il contesto che ha motivato la navigazione.

La schermata denominata \emph{Valutation Success Criteria} raccoglie i sedici criteri definiti nella fonte di ricerca. I filtri sono organizzati in due gruppi: il primo segue i Principi, mentre il secondo espone le Linee Guida per le quali sono disponibili criteri. Ciascuna voce mostra il codice sintetico, il nome esteso della Linea Guida e il numero di criteri associati. Questa ridondanza controllata evita che sigle quali G1 o G4 debbano essere ricordate senza supporto. I criteri sono presentati tramite card espandibili: nello stato iniziale sono visibili codice e formulazione principale; l'attivazione della freccia interna mostra la descrizione e i riferimenti al Principio e alla Linea Guida. È stata adottata una visualizzazione progressiva perché la presenza simultanea delle sedici descrizioni avrebbe reso la pagina eccessivamente densa, riducendo la possibilità di confrontare rapidamente i diversi output di valutazione.

La Design Pattern Library utilizza una struttura analoga, ma adeguata alla maggiore quantità di informazioni contenuta in ciascuna scheda. Le card sono state ridimensionate e distribuite su una griglia che può raggiungere tre colonne. Nello spazio disponibile vengono mostrati codice, titolo, problema, suggerimento di design e Linee Guida collegate. La porzione inferiore segnala inoltre la presenza di un esempio, destinato a essere accompagnato dalla relativa figura esplicativa. La sintesi non sostituisce la scheda completa: il collegamento \emph{Explore} conduce alla pagina del singolo pattern, nella quale sono esposti Problem, Design Suggestion, Solution, Context and Motivations ed Examples. La freccia dell'intestazione riconduce sempre alla Design Pattern Library, mentre le Linee Guida mostrate a lato sono selezionabili e riportano alla rispettiva card nella loro libreria.

Il comportamento dei filtri rispecchia la struttura reticolare del framework. Una card può appartenere a più categorie e rimane visibile quando almeno una delle associazioni coincide con il filtro scelto. Lo stato selezionato viene comunicato sia visivamente sia mediante attributi semantici, mentre un messaggio dedicato gestisce il caso in cui non esistano risultati. Parte dello stato può essere espresso anche nell'indirizzo della pagina. Per esempio, il passaggio da una Linea Guida ai relativi Design Pattern aggiunge il codice della guideline all'URL; la pagina di destinazione interpreta tale parametro e inizializza il filtro corrispondente. Questo accorgimento permette di conservare il significato del collegamento, condividere una vista già filtrata e utilizzare correttamente i comandi avanti e indietro del browser.

Accanto ai percorsi trasversali sono stati mantenuti comandi di ritorno coerenti con la gerarchia. Nelle quattro schermate principali la freccia conduce alla homepage; nelle pagine dei Principi individuali conduce alla Library of Principles; nelle schede dei singoli Design Pattern conduce alla Design Pattern Library. La decisione evita che lo stesso simbolo produca risultati diversi in base al punto dal quale la pagina è stata aperta. Il logo collocato nella barra laterale costituisce infine un secondo accesso costante alla homepage. Navigazione globale e navigazione gerarchica restano così distinte: la prima consente di cambiare area, la seconda permette di risalire dal dettaglio alla raccolta di appartenenza.

\section{Tecniche implementative e di sviluppo}

FrameSAI è stata sviluppata come applicazione server-side basata su Laravel 12 e PHP 8.2. L'adozione del framework ha permesso di organizzare il progetto secondo il pattern Model--View--Controller, separando la risoluzione delle richieste, la preparazione dei contenuti e la loro rappresentazione. La versione iniziale disponeva già di rotte, controller e viste Blade, ma ciascuna schermata conteneva una parte consistente della struttura HTML comune e definiva localmente i dati da mostrare. Questa impostazione consentiva di riprodurre rapidamente il prototipo, ma produceva duplicazione: la barra laterale, l'intestazione, il caricamento degli asset e numerose regole visuali erano ripetuti in ogni file.

La fase di consolidamento ha introdotto un layout Blade condiviso. Il file di layout contiene la struttura generale del documento, i metadati, il caricamento degli asset e le aree nelle quali ogni pagina inserisce titolo, descrizione e contenuto. La navigazione principale e l'intestazione sono state isolate in componenti riutilizzabili. Il primo riceve l'indicazione della sezione attiva e costruisce in modo uniforme le cinque destinazioni della piattaforma; il secondo riceve titolo, sottotitolo, icona e destinazione del comando di ritorno. Una modifica alla sidebar o alla dimensione delle icone viene quindi applicata a tutte le schermate, evitando divergenze generate da correzioni locali.

I controller costituiscono il punto di raccordo tra le richieste e le viste. \texttt{LibraryPrinciplesController}, \texttt{GuidelinesController}, \texttt{SuccessCriteriaController} e \texttt{DesignPatternController} preparano le collezioni necessarie alle rispettive librerie. I controller delle pagine di dettaglio ricevono invece un identificativo dalla rotta, verificano che esso corrisponda a una risorsa esistente e, in caso contrario, restituiscono una risposta 404. Questa gestione impedisce la produzione di pagine formalmente valide ma prive di contenuto. Gli indirizzi sono stati formulati in modo leggibile, come \texttt{/principles}, \texttt{/success-criteria} e \texttt{/design-patterns/4-1}; i precedenti indirizzi delle librerie vengono reindirizzati verso le nuove destinazioni per non interrompere eventuali collegamenti già utilizzati.

Nella versione corrente il catalogo del framework è raccolto in un file di configurazione dedicato. Principi, Linee Guida, Criteri di Successo e Design Pattern sono descritti tramite strutture associative che includono codici, testi e collegamenti. Questa soluzione rappresenta un passaggio intermedio tra il contenuto incorporato nelle viste e un livello di persistenza completo. Ha consentito di centralizzare la fonte utilizzata dall'applicazione e di calcolare dinamicamente i conteggi mostrati nei filtri, senza introdurre un database prima di stabilizzare il modello delle relazioni. La piattaforma non offre ancora operazioni di inserimento o modifica da interfaccia, autenticazione o salvataggio di progetti personali; tali funzioni richiederebbero modelli persistenti e costituiscono un'estensione successiva, non una caratteristica simulata della versione descritta.

Le viste sono realizzate con Blade e HTML semantico, mentre Tailwind CSS 4 gestisce la presentazione. L'utilizzo di classi responsive permette di modificare griglie, spaziature e orientamento dei componenti in corrispondenza della larghezza disponibile. La prima trasposizione del prototipo imponeva larghezze minime comprese tra 1180 e 1280 pixel, condizione che generava uno scorrimento orizzontale sui dispositivi più piccoli. Tali vincoli sono stati rimossi: le card passano da una a più colonne, i filtri vengono collocati sopra il contenuto quando lo spazio laterale è insufficiente e la navigazione verticale assume una disposizione orizzontale. È stata mantenuta una larghezza massima per le aree testuali, così da non sacrificare la leggibilità sui monitor ampi.

Le interazioni circoscritte, quali filtri e pannelli espandibili, sono state sviluppate in JavaScript senza introdurre un framework client aggiuntivo. Per i filtri, lo script legge le associazioni esposte negli attributi \texttt{data-*}, confronta il valore selezionato e modifica la visibilità delle card. Nei Success Criteria, il comando interno aggiorna contemporaneamente la sezione mostrata, l'orientamento dell'icona e il valore di \texttt{aria-expanded}. La scelta di JavaScript nativo è proporzionata alla complessità della versione corrente: l'interfaccia non possiede uno stato applicativo globale né flussi di modifica che giustifichino una seconda architettura frontend.

Vite viene utilizzato per compilare e ottimizzare CSS e JavaScript. La separazione tra file sorgente e asset generati consente di utilizzare le classi di Tailwind durante lo sviluppo e di produrre un insieme compatto di risorse per l'esecuzione. La pipeline rimane tuttavia subordinata alla struttura Laravel: il template engine genera il documento richiesto, mentre gli asset compilati ne determinano stile e comportamento nel browser.

Particolare attenzione è stata riservata alla semantica dell'interfaccia. Le regioni principali utilizzano elementi quali \texttt{main}, \texttt{nav}, \texttt{aside}, \texttt{section} e \texttt{article}; la voce di navigazione corrente espone \texttt{aria-current}, i filtri comunicano la selezione tramite \texttt{aria-pressed} e i pannelli espandibili utilizzano \texttt{aria-expanded}. Le icone decorative possiedono testo alternativo vuoto e non intercettano gli eventi del puntatore. Quest'ultimo intervento è risultato necessario dopo l'ingrandimento delle immagini viola collocate vicino ai titoli: la trasformazione visuale poteva sovrapporre il riquadro dell'immagine alla freccia di ritorno, rendendola apparentemente non funzionante. La freccia è stata quindi portata su un livello interattivo superiore e le immagini decorative sono state escluse dalla gestione del puntatore.

La verifica tecnica affianca lo sviluppo. La suite di test controlla la risposta delle pagine pubbliche, tutte le schede dei Design Pattern, le risorse non esistenti, i reindirizzamenti e le destinazioni dei comandi di ritorno. Nello stato descritto, sei test raccolgono settantotto asserzioni. A questa verifica automatica si aggiunge il controllo nel browser a differenti risoluzioni: sono stati osservati l'assenza di overflow orizzontale, la permanenza della sidebar durante lo scorrimento, il funzionamento dei filtri e l'effettiva navigazione prodotta dalle frecce. I test non sostituiscono la valutazione con utenti, ma riducono il rischio che una modifica strutturale interrompa i percorsi necessari alla successiva fase sperimentale.

\section{Funzionalità: output e motivazioni}

L'output principale della piattaforma è una rappresentazione consultabile del framework FrameSAI. Il valore dell'implementazione non consiste soltanto nell'aver trasferito i contenuti dei documenti in un ambiente web, ma nell'aver reso esplicite le relazioni che nei materiali originari richiedevano una lettura trasversale. I codici P, G, SC e DP sono conservati perché consentono un riferimento puntuale, mentre titoli, descrizioni e conteggi ne riducono l'astrazione. La piattaforma combina quindi due modalità di lettura: una sintetica, adatta a individuare categorie e corrispondenze, e una analitica, destinata all'approfondimento delle singole risorse.

La Library of Principles produce come output una visione compatta dei requisiti generali della simbiosi. La scelta delle card deriva dalla necessità di presentare elementi omogenei senza introdurre una graduatoria: i quattro Principi condividono lo stesso livello e differiscono attraverso colore, codice e contenuto. La pagina individuale amplia l'informazione mediante una descrizione e un esempio, offrendo un passaggio dalla definizione astratta a una possibile manifestazione nell'interfaccia. Il collegamento alle altre librerie segnala che un Principio non è verificato direttamente, ma viene progressivamente specializzato da Linee Guida e Criteri.

La libreria delle Linee Guida restituisce invece un insieme di prescrizioni e la rete delle relative appartenenze. La presenza dei conteggi non è stata utilizzata come indicatore di importanza: una Linea Guida con più pattern non è necessariamente prioritaria rispetto a una con un solo collegamento. Il numero comunica piuttosto l'estensione del materiale disponibile e permette di anticipare l'esito della selezione. Le card G5, G8 o G9 possono, per esempio, non disporre dello stesso numero di output implementativi di G2 o G3; il sistema mantiene visibile questa asimmetria anziché colmarla con associazioni non presenti nelle fonti.

I Success Criteria costituiscono l'output valutativo della piattaforma. La loro formulazione è presentata come condizione osservabile, non come punteggio già calcolato. La versione corrente non assegna automaticamente una valutazione a un sistema e non produce un esito di conformità: rende disponibili i criteri attraverso i quali tale valutazione può essere condotta. Questa distinzione è rilevante perché evita di trasformare istruzioni qualitative in una misura numerica priva di una procedura validata. L'espansione della card permette di accedere alla spiegazione del criterio e ai riferimenti necessari, lasciando al valutatore il compito di raccogliere evidenze nel sistema analizzato.

La Design Pattern Library produce infine l'output più vicino all'attività progettuale. Il Problema definisce la condizione che richiede un intervento; il Design Suggestion individua l'orientamento della soluzione; la Solution articola le operazioni possibili; Context and Motivations chiarisce quando il pattern acquista rilevanza; Examples collega la prescrizione a un caso o a una tecnica esistente. Questa struttura è stata mantenuta anche quando il numero di punti nella soluzione varia, evitando di comprimere ogni pattern in una forma identica soltanto per uniformità visuale. Nelle card della libreria viene mostrato un sottoinsieme delle informazioni, mentre la scheda individuale conserva il testo completo disponibile.

La scelta di non mostrare collegamenti ai Success Criteria nelle pagine individuali dei pattern risponde allo stato effettivo del materiale. Le associazioni tra pattern e guideline sono esplicite e utilizzate per la navigazione; un collegamento diretto e definitivo tra ogni pattern e specifici criteri richiederebbe invece una validazione ulteriore. La piattaforma evita pertanto di presentarlo come dato acquisito. Nelle card viene indicato \emph{1 example}, riferendosi alla figura esplicativa prevista nella sezione Examples. Tale spazio potrà ospitare progressivamente le immagini contenute nel documento dei pattern senza modificare la struttura della libreria.

I filtri rappresentano una funzione di orientamento, non soltanto un modo per ridurre il numero di card. Nella Design Pattern Library consentono di partire da una Linea Guida e osservare le soluzioni che la concretizzano; nei Success Criteria permettono di cambiare prospettiva tra Principio e Linea Guida; nella schermata delle Linee Guida evidenziano le sovrapposizioni tra i quattro Principi. Lo stesso contenuto può quindi essere interrogato secondo la domanda dell'utente: comprendere una proprietà generale, individuare un obbligo, cercare una condizione valutabile oppure reperire una soluzione di interaction design.

La persistenza della dashboard laterale sostiene questi passaggi rendendo visibile l'intera architettura anche durante la lettura di pagine lunghe. La scelta di mantenerla fissa non serve a incentivare cambi frequenti di sezione, ma a evitare che l'utente percepisca le pagine di dettaglio come documenti isolati. Allo stesso modo, i colori identificativi non hanno una funzione puramente decorativa: producono continuità tra la card del Principio, il filtro e il badge visualizzato nelle risorse correlate. Il codice testuale affiancato al colore mantiene l'associazione interpretabile anche quando le tonalità non sono distinguibili.

Nel complesso, FrameSAI restituisce un artefatto dimostrativo coerente con gli obiettivi della fase di progettazione. La piattaforma permette di attraversare il framework, confrontarne gli elementi e raggiungere il dettaglio senza perdere il contesto. Non costituisce ancora un sistema di assessment automatico né un ambiente collaborativo nel quale salvare progetti e risultati; costituisce però la base tecnica necessaria per entrambe le evoluzioni. La separazione dei contenuti, la struttura delle rotte e i componenti condivisi consentono infatti di sostituire in seguito il catalogo configurato con modelli persistenti, mantenendo invariati i principali percorsi di consultazione. Il contributo dell'implementazione risiede pertanto nel rendere la struttura teorica osservabile e utilizzabile, predisponendola alla successiva valutazione con gli utenti target.
